\section{Constant profiles}
\label{sec:constant-profiles}

On \(\mathbb C[t]\), the sphere functional \(L\) is the interval average
\(J(f)=\frac12\int_{-1}^1 f(t)\,dt\).

\begin{lemma}[Interval cubics]
\label{lem:interval-cubic}
If \(u\in\mathbb C[t]\), \(\deg u\le3\), and
\(J(u^m)=0\) for \(1\le m\le4\), then \(u=0\).
\end{lemma}

\begin{proof}
The first equation puts
\(u=h t^3+v(t^2-1/3)+ct\). Write \(M_j=J(u^j)\). Direct integration gives
\begin{align}
 M_2&=\frac{105c^2+126ch+45h^2+28v^2}{315},\notag\\
 M_3\bmod M_2&=\frac{8v(56c^2+105ch+45h^2)}{2205},\notag\\
 M_4\bmod M_2&=-\frac{16R(c,h)}{735735},
 \label{eq:interval-cubic-remainders}
\end{align}
where remainders are taken as polynomials in \(v\), and
\[
 R(c,h)=28028c^4+84084c^3h+93275c^2h^2+46410ch^3+8955h^4.
\]
If \(h=0\), division of \(M_3\) by \(M_2\), now in the variable
\(c\), has remainder \(-256v^3/4725\). Thus \(v=c=0\).

Suppose \(h\ne0\). Scale \(h=1\), rename \(c=z\), and put
\[
 P=56z^2+105z+45,\qquad B=105z^2+126z+45,\qquad R=R(z,1).
\]
If \(v\ne0\), the equations force \(P=R=0\). But
\[
 R\bmod P=-\frac{315}{128}(299z+159),\qquad
 P(-159/299)=\frac{446976}{89401}\ne0.
\]
If \(v=0\), they force \(B=R=0\), whereas
\[
 R\bmod B=-\frac{84}{125}(117z-110),\qquad
 B(110/117)=\frac{1169375}{4563}\ne0.
\]
Both alternatives are impossible.
\end{proof}

The geometric reduction supplies either a phase independent of \(y\),
or a constant profile in which the phase-zero term \(t^3\) is absent.
The first case is already settled by Lemma~\ref{lem:interval-cubic}:
if \(p\in\mathbb C[x,t]\), phase averaging leaves its coefficient of
\(x^0\); purity makes that interval cubic zero. All remaining terms
have strictly positive circular weight.

\begin{proposition}[Constant-profile normal forms]
\label{prop:finite-constant-profile}
Let \(a\ne0\) and
\begin{equation}
 p=ay+(v_0+v_1x)t^2+(c+dx+ex^2)t+\sum_{j=0}^3r_jx^j.
 \label{eq:constant-profile-family}
\end{equation}
If \(L(p^m)=0\) for \(1\le m\le8\), then \(v_0=v_1=e=0\) and
\begin{equation}
 \sum_{j=0}^3r_jx^j=-\frac{x(c+dx)^2}{2a}.
 \label{eq:constant-profile-shear}
\end{equation}
\end{proposition}

\begin{proof}
Divide by \(a\), rename coefficients, and assume \(a=1\).
The orthogonal null rotation
\begin{equation}
 x=x',\qquad t=t'+kx',\qquad y=y'-kt'-k^2x'/2
 \label{eq:constant-profile-null-rotation}
\end{equation}
preserves \eqref{eq:constant-profile-family}, leaves \(v_0,v_1\)
unchanged, and replaces \(c\) by \(c-k\). Take \(k=c\).

Here is a closed formula for the finite coefficient calculation. For
\[
 p=y+\sum_{\nu=1}^N a_\nu x^{i_\nu}t^{j_\nu},\qquad i_\nu\ge0,
\]
let \(\alpha\in\mathbb N^N\),
\(k=\sum i_\nu\alpha_\nu\), and \(j=\sum j_\nu\alpha_\nu\).
Phase balance forces exactly \(k\) factors of \(y\), so
\begin{equation}
 L(p^m)=
 \sum_{\substack{\sum(1+i_\nu)\alpha_\nu=m\\j\ \mathrm{even}}}
 \frac{m!(j-1)!!}{(2k+j+1)!!\prod_\nu\alpha_\nu!}
 \prod_\nu a_\nu^{\alpha_\nu}.
 \label{eq:constant-profile-coefficients}
\end{equation}
The factorial for the \(y\)-factors cancels the \(k!\) in
\eqref{eq:moments}. Thus this formula computes every
coefficient over \(\mathbb Q\), for all parameter values at once.

For \(a=1,c=0\), the first four equations solve successively for
\(r_0,r_1,r_2,r_3\). Their respective coefficients are the nonzero
constants \(1,2/3,2/5,8/35\), and the solutions are
\begin{align}
 r_0&=-v_0/3,&
 r_1&=-(2v_0^2+3v_1)/15,\notag\\
 r_2&=-4v_0(31v_0^2+54v_1)/945,\notag\\
 r_3&=-d^2/2-2312v_0^4/14175-188v_0^2v_1/525-8v_1^2/75.
 \label{eq:constant-profile-triangular}
\end{align}
After substitution, if \(v_0=0\), the eighth equation is
\(-65536v_1^4/41023125=0\). Then the sixth equation becomes
\(8e^2/231=0\). Hence \(v_1=e=0\), and
\eqref{eq:constant-profile-triangular} leaves only \(r_3=-d^2/2\).

Suppose \(v_0\ne0\), and set
\[
 b=v_1/v_0^2,\qquad D=d^2/v_0^4,\qquad E=de/v_0^5,\qquad F=e^2/v_0^6.
\]
Let \(N_i\), \(5\le i\le8\), be the primitive integer numerators of
the four residual moments, with positive coefficient of \(v_0^i\).
The equations \(N_i/v_0^i=0\) are linear in \(D,E,F\). Their augmented
matrix, in the column order \(D,E,F,1\), is
\begin{equation}
 M(b)=\begin{pmatrix}
 155925&155925&0&B_5\\
 77395500-116093250b&77395500&-58046625&B_6\\
 226800-297675b&226800+127575b&-212625&B_7\\
 4819500-216877500b&4819500-65063250b&-75907125&B_8
 \end{pmatrix},
 \label{eq:constant-profile-matrix}
\end{equation}
where
\begin{align*}
 B_5&=35504+94992b+52704b^2,\\
 B_6&=30907712+63343584b+3500496b^2-12573792b^3,\\
 B_7&=53504+17328b-211968b^2-108864b^3,\\
 B_8&=7986496-27544704b-90351504b^2\\
    &\qquad{}-16632864b^3+9144576b^4.
\end{align*}
Its coefficient columns satisfy \(C_D-C_E-2bC_F=0\). Consistency
therefore forces every three-by-three augmented minor to vanish.
Columns \(D,F,1\), with the first three rows and with rows one, two,
four, give, up to nonzero rational factors,
\begin{align*}
 f(b)&=2485161b^3+6647562b^2+4908168b+1097984,\\
 g(b)&=10216206b^4+24356619b^3-25372278b^2-41158800b-11711872.
\end{align*}
In \(\mathbb F_{19}[b]\), the reductions are
\[
 \bar f=-b^3-6b^2-7b-7,\qquad
 \bar g=b^4+6b^3-b^2+7b-6.
\]
Starting with \((\bar g,\bar f)\), the successive Euclidean remainders are
\[
 -8b^2-6,\qquad 8b+7,\qquad -5.
\]
Hence the reductions are coprime. Both leading coefficients remain
nonzero modulo \(19\). By Gauss's lemma, a nonconstant common factor
over \(\mathbb Q\) could be chosen primitive over \(\mathbb Z\); its
leading coefficient would be a unit modulo \(19\), and its reduction
would remain a nonconstant common factor. This is impossible.
This excludes \(v_0\ne0\), without dividing by \(b\)
or assuming any further coefficient is nonzero. Undoing
\eqref{eq:constant-profile-null-rotation} and the scalar normalization
gives \eqref{eq:constant-profile-shear}.

\end{proof}

